\begin{figure}[ht]
\centering
\begin{tikzpicture}[x=0.85cm, y=0.4cm, font=\small]

% ----- Panel A: centred process (top) -----
\begin{scope}[yshift=14cm]
\draw[->, thick] (-0.2, 0) -- (12, 0)
  node[right, font=\scriptsize] {$L - 75 \,(\si{\milli\meter})$};
\draw[->, thick] (0, -0.1) -- (0, 11);

% x-axis ticks (centred at 0, span +/- 0.10 mm scaled to 0..12)
\foreach \x/\lbl in {1/-0.10, 3/-0.05, 6/0.00, 9/0.05, 11/0.10} {
  \draw (\x, 0) -- (\x, -0.3) node[below, font=\scriptsize] {\lbl};
}

% Centred-process histogram (RSS sd ~ 0.029 mm; bins around x=6)
\foreach \x/\h in {
  2/0.3, 2.5/0.6, 3/1.2, 3.5/2.3, 4/3.8,
  4.5/5.8, 5/7.8, 5.5/9.2, 6/9.8, 6.5/9.2,
  7/7.8, 7.5/5.8, 8/3.8, 8.5/2.3, 9/1.2,
  9.5/0.6, 10/0.3} {
  \fill[engblue!30] (\x, 0) rectangle (\x + 0.5, \h);
  \draw[engblue!50] (\x, 0) rectangle (\x + 0.5, \h);
}
% Gaussian overlay (linear-propagation prediction)
\draw[thick, engred] plot[domain=2:10, samples=80]
  (\x, {9.8*exp(-((\x-6)*(\x-6))/(2*1.7*1.7))});

\node[anchor=south, font=\small] at (6, 11) {(a) centred process};
\node[engblue, anchor=west, font=\scriptsize] at (8.0, 8.5) {MC, truncated};
\node[engred, anchor=west, font=\scriptsize] at (8.0, 7.5) {RSS Gaussian};
\end{scope}

% ----- Panel B: shifted process (bottom) -----
\begin{scope}
\draw[->, thick] (-0.2, 0) -- (12, 0)
  node[right, font=\scriptsize] {$L - 75 \,(\si{\milli\meter})$};
\draw[->, thick] (0, -0.1) -- (0, 11);

\foreach \x/\lbl in {1/-0.10, 3/-0.05, 6/0.00, 9/0.05, 11/0.10} {
  \draw (\x, 0) -- (\x, -0.3) node[below, font=\scriptsize] {\lbl};
}

% Shifted-process histogram: mean at +0.025 mm (between 6 and 9 on plot),
% asymmetric upper tail truncated at spec limit.
\foreach \x/\h in {
  3/0.2, 3.5/0.5, 4/1.0, 4.5/2.0, 5/3.4,
  5.5/5.2, 6/7.2, 6.5/8.6, 7/9.6, 7.5/9.8,
  8/9.2, 8.5/7.6, 9/5.0, 9.5/2.0, 10/0.5} {
  \fill[engamber!40] (\x, 0) rectangle (\x + 0.5, \h);
  \draw[engamber!70] (\x, 0) rectangle (\x + 0.5, \h);
}
% RSS Gaussian centred at nominal (x=6) for comparison
\draw[thick, engred, dashed] plot[domain=2:10, samples=80]
  (\x, {9.8*exp(-((\x-6)*(\x-6))/(2*1.7*1.7))});

\node[anchor=south, font=\small] at (6, 11) {(b) shifted process};
\node[engamber, anchor=west, font=\scriptsize] at (8.0, 8.5) {MC, truncated};
\node[engred, anchor=west, font=\scriptsize] at (8.0, 7.5) {RSS Gaussian};
\end{scope}

\end{tikzpicture}
\caption[Tolerance stack: linear vs Monte Carlo under truncation]{
Empirical distributions of assembly length $L = a + b + c$ for two
manufacturing scenarios. Panel (a): centred process with all three
part means at nominal, manufacturing distribution Gaussian truncated
at the specification limit. The truncation barely bites and the
Monte Carlo distribution agrees with the linear-propagation RSS
Gaussian to plotting accuracy. Panel (b): shifted process with part
1's manufacturing mean offset by $+t/2$ inside the upper
specification. The truncation removes a substantial fraction of
the upper tail, the empirical distribution is visibly asymmetric,
and the assembly mean lies above the nominal of $75\,\si{\milli
\meter}$. The RSS Gaussian centred at nominal (dashed) understates
the assembly mean and overstates the upper-tail width.}
\label{fig:vol01:ch04:tolerance-stack}
\end{figure}
